% BHU PYQ -- Auto-generated & error-corrected
\documentclass[11pt,a4paper]{article}

\usepackage[a4paper,top=2.5cm,bottom=2.5cm,left=2.8cm,right=2.8cm,headheight=14pt]{geometry}
\usepackage{amsmath,amssymb}
\usepackage[T1]{fontenc}
\usepackage[utf8]{inputenc}
\usepackage{lmodern}
\usepackage{microtype}
\usepackage{fancyhdr}
\usepackage{enumitem}
\usepackage{hyperref}
\usepackage{lastpage}
\usepackage{parskip}

\hypersetup{colorlinks=true,urlcolor=black,linkcolor=black,
  pdftitle={STA/STB-502: Statistical Inference and Decision Theory -- BHU PYQ},pdfauthor={Banaras Hindu University}}

\pagestyle{fancy}\fancyhf{}
\renewcommand{\headrulewidth}{0.4pt}\renewcommand{\footrulewidth}{0.4pt}
\fancyhead[L]{\small\textit{Banaras Hindu University}}
\fancyhead[R]{\small\textit{STA/STB-502: Statistical Inference \& Decision Theory}}
\fancyfoot[C]{\small Page \thepage\ of \pageref{LastPage}}

\newlist{parts}{enumerate}{2}
\setlist[parts,1]{label=(\alph*),leftmargin=*,itemsep=5pt,topsep=3pt}
\setlist[parts,2]{label=(\roman*),leftmargin=*,itemsep=3pt,topsep=2pt}
\newcommand{\pts}[1]{\hfill\textit{[#1]}}

\begin{document}

\begin{center}
  {\large\bfseries BANARAS HINDU UNIVERSITY}\\[2pt]
  {\normalsize B.A./B.Sc. (Hons.) Semester V Examination, 2016-17}\\[6pt]
  \rule{0.8\linewidth}{0.5pt}\\[6pt]
  {\large\bfseries Paper No. STA/STB-502: Statistical Inference and Decision Theory}\\[4pt]
  {\normalsize Time: 3 Hours\hfill Maximum Marks: 70}
\end{center}
\rule{\linewidth}{0.4pt}
\smallskip
\noindent\textit{Instructions: Answer any \textbf{five} questions. All questions carry equal marks. Symbols have their usual meanings.}
\bigskip

\medskip\noindent\textbf{Question 1.}
\begin{parts}
  \item Define consistency of an estimator. Let $X_1, X_2, \ldots, X_n$ be a random sample from the probability density function:
  \[
    f(x, \theta) = e^{-(x-\theta)}, \quad x \geq \theta, \quad \theta > 0.
  \]
  Verify whether $T_n = X_{(1)}$ (the first-order statistic) is consistent for $\theta$. \pts{7}
  \item If $T_1$ is the minimum variance unbiased estimator (MVUE) of $\psi(\theta)$, and $T_2$ is any other unbiased estimator of $\psi(\theta)$ with efficiency $e = e_0$, then show that the correlation coefficient between $T_1$ and $T_2$ is given by:
  \[
    \rho = \sqrt{e_0}.
  \]\pts{7}
\end{parts}

\medskip\rule{\linewidth}{0.2pt}

\medskip\noindent\textbf{Question 2.}
\begin{parts}
  \item Define a sufficient statistic. Let $X_1, X_2, \ldots, X_n$ be a random sample of size $n$ from a Poisson distribution with parameter $\theta$. Verify whether (i) $X_1 + X_2$ and (ii) $X_1 + 2X_2$ are sufficient statistics for $\theta$. \pts{7}
  \item Let $X_1, X_2, \ldots, X_n$ be a random sample of size $n$ from $U(0, \theta)$. Find the minimum variance unbiased estimator of $\theta$. \pts{7}
\end{parts}

\medskip\rule{\linewidth}{0.2pt}

\medskip\noindent\textbf{Question 3.}
\begin{parts}
  \item State and prove the Rao--Blackwell theorem. \pts{7}
  \item What do you mean by a minimum variance bound estimator? Let $X_1, X_2, \ldots, X_n$ be a random sample from:
  \[
    f(x, \theta) = \theta e^{-\theta x}, \quad x \geq 0, \quad \theta > 0.
  \]
  Obtain the minimum variance bound estimator of $\theta$. \pts{7}
\end{parts}

\medskip\rule{\linewidth}{0.2pt}

\medskip\noindent\textbf{Question 4.}
\begin{parts}
  \item Explain the method of maximum likelihood estimation. Obtain the MLE of $\theta$ based on a random sample of size $n$ from the following distributions:
  \begin{parts}
    \item $f(x, \theta) = \theta x^{\theta-1}$, \quad $0 < x < 1$, \quad $\theta > 0$
    \item $f(x, \theta) = \theta(1-x)^{\theta-1}$, \quad $0 < x < 1$, \quad $\theta > 0$
  \end{parts}\pts{7}
  \item Explain the method of moments for estimating parameters. Let $X_1, X_2, \ldots, X_n$ be a random sample of size $n$ from a Binomial distribution $B(n, p)$ with unknown parameters $n$ and $p$. Obtain the estimators of the parameters by the method of moments. \pts{7}
\end{parts}

\medskip\rule{\linewidth}{0.2pt}

\medskip\noindent\textbf{Question 5.}
\begin{parts}
  \item What is a confidence interval? What is meant by confidence coefficient? Explain the method of construction of confidence intervals. \pts{7}
  \item Obtain the $100(1-\alpha)\%$ confidence interval for $\sigma^2$ of the normal distribution $N(\mu, \sigma^2)$ when:
  \begin{parts}
    \item $\mu$ is known,
    \item $\mu$ is unknown.
  \end{parts}\pts{7}
\end{parts}

\medskip\rule{\linewidth}{0.2pt}

\medskip\noindent\textbf{Question 6.}
\begin{parts}
  \item Explain the following terms:
  \begin{parts}
    \item Simple and composite hypothesis
    \item Most powerful test
    \item Size of a test and level of significance
  \end{parts}\pts{7}
  \item State the Neyman--Pearson lemma. Obtain the most powerful test for testing $H_0\colon \theta = \theta_0$ against $H_1\colon \theta = \theta_1\;(\theta_1 > \theta_0)$ on the basis of a sample of size $n$ from an exponential distribution with pdf:
  \[
    f(x, \theta) = \theta e^{-\theta x}, \quad x > 0, \quad \theta > 0.
  \]\pts{7}
\end{parts}

\medskip\rule{\linewidth}{0.2pt}

\medskip\noindent\textbf{Question 7.}

Describe the procedure of the likelihood ratio test. Obtain the likelihood ratio test for testing $H_0\colon \sigma^2 = \sigma_0^2$ against $H_1\colon \sigma^2 \neq \sigma_0^2$ on the basis of a random sample of size $n$ from $N(\mu, \sigma^2)$ with unknown mean $\mu$. \pts{14}

\medskip\rule{\linewidth}{0.2pt}

\medskip\noindent\textbf{Question 8.}

Explain the loss function and the risk function. Let $X_1, X_2, \ldots, X_n$ be a random sample from $N(\mu, 1)$. Under the squared error loss function, obtain the Bayes estimator of $\mu$ when the prior distribution for $\mu$ is $N(\mu_0, \tau^2)$. \pts{14}

\vfill
\rule{\linewidth}{0.4pt}
\begin{center}\small
  \href{https://bhu-exam.vercel.app/}{Click here to visit our website}\\[4pt]
  \href{https://me-aryan.vercel.app/}{Click here to visit our contributor}\\[4pt]
  \href{https://quashan-me.vercel.app/}{Click here to visit our contributor}
\end{center}

\end{document}